\section{Exact Density Modeling Through Invertible Maps}

A third major route to generative modeling is based on exact transformation of probability distributions.
This leads to \emph{normalizing flows}: models that transport a simple base distribution into a more complex data distribution through an invertible map.

\medskip

Compared with adversarial models, flows retain exact density evaluation.
Compared with latent-variable models such as VAEs, they avoid variational approximation in the likelihood itself.
Their main burden is architectural: the map must be expressive, invertible, and accompanied by a tractable Jacobian determinant.

\subsection*{Invertibility as a modeling strategy}

Let $z\in\mathbb{R}^d$ be drawn from a simple base density $p_Z$, for example a standard Gaussian, and let
\[
x = f_\theta(z),
\]
where $f_\theta:\mathbb{R}^d\to\mathbb{R}^d$ is bijective and differentiable with differentiable inverse.
Then sampling is straightforward:
\[
z\sim p_Z, \qquad x=f_\theta(z).
\]

Because $f_\theta$ is invertible, one may also map a data point $x$ back to latent coordinates
\[
z=f_\theta^{-1}(x),
\]
and evaluate its density exactly by accounting for local volume distortion.
This is the central mathematical idea behind flows.

\subsection*{The finite-dimensional change-of-variables theorem}

\begin{theorem}[Change of variables in finite dimensions]
Let $f:U\to V$ be a $C^1$ bijection between open sets $U,V\subset \mathbb{R}^d$, and assume that the Jacobian matrix $J_f(z)$ is nonsingular for all $z\in U$.
If $Z$ has density $p_Z$ on $U$ and $X=f(Z)$, then $X$ has density
\[
p_X(x)=p_Z\bigl(f^{-1}(x)\bigr)\left|\det J_{f^{-1}}(x)\right|,
\qquad x\in V.
\]
Equivalently, writing $x=f(z)$,
\[
p_X\bigl(f(z)\bigr)=p_Z(z)\left|\det J_f(z)\right|^{-1}.
\]
\end{theorem}

\begin{proof}[Proof sketch]
For small regions in latent space, a differentiable map behaves approximately like its Jacobian matrix.
Thus a small volume element $dz$ near $z$ is mapped to a volume element near $x=f(z)$ whose size is approximately
\[
\left|\det J_f(z)\right|dz.
\]
Probability mass is preserved under deterministic transformation, so
\[
p_Z(z)\,dz \approx p_X(x)\,dx = p_X(f(z))\left|\det J_f(z)\right|dz.
\]
Dividing by $dz$ yields
\[
p_X(f(z)) = p_Z(z)\left|\det J_f(z)\right|^{-1}.
\]
Replacing $z$ by $f^{-1}(x)$ gives the first formula.
\end{proof}

\paragraph{Interpretation.}
The determinant measures how much the map expands or contracts local volume.
If $|\det J_f(z)|>1$, the transformation spreads probability mass over a larger region and the resulting density decreases.
If $|\det J_f(z)|<1$, probability mass is concentrated and the density increases.
Thus exact density modeling is possible whenever the transport map and its local volume change are tractable.

\subsection*{Likelihood under an invertible model}

Given data $x\in\mathbb{R}^d$ and an invertible model $x=f_\theta(z)$ with base density $p_Z$, the log-likelihood is
\[
\log p_\theta(x)
=
\log p_Z\bigl(f_\theta^{-1}(x)\bigr)
+
\log\left|\det J_{f_\theta^{-1}}(x)\right|.
\]
Equivalently, if $z=f_\theta^{-1}(x)$, then
\[
\log p_\theta(x)
=
\log p_Z(z)-\log\left|\det J_{f_\theta}(z)\right|.
\]

This is one reason flows are mathematically elegant:
\begin{itemize}
    \item sampling is exact,
    \item density evaluation is exact,
    \item training can proceed by direct maximum likelihood.
\end{itemize}

\subsection*{Composition of invertible layers}

A practical flow is usually a composition of simpler invertible maps,
\[
f_\theta = f_L\circ f_{L-1}\circ \cdots \circ f_1.
\]
If we define intermediate states
\[
h_0=z, \qquad h_\ell=f_\ell(h_{\ell-1}), \qquad h_L=x,
\]
then the chain rule gives the total Jacobian
\[
J_{f_\theta}(z)=J_{f_L}(h_{L-1})\cdots J_{f_1}(h_0).
\]

\begin{proposition}[Composition formula for log-determinants]
For the composition above,
\[
\log \left|\det J_{f_\theta}(z)\right|
=
\sum_{\ell=1}^L \log\left|\det J_{f_\ell}(h_{\ell-1})\right|.
\]
\end{proposition}

\begin{proof}
The determinant of a product of square matrices equals the product of determinants:
\[
\det(AB)=\det(A)\det(B).
\]
Applying this repeatedly to the chain-rule factorization and then taking logarithms gives the result.
\end{proof}

\paragraph{Why this matters.}
The overall transformation may be highly nonlinear, yet exact density evaluation reduces to summing layerwise contributions.
Thus the engineering problem becomes: design invertible layers for which the inverse and the Jacobian determinant are both cheap.

\subsection*{Coupling layers and tractable Jacobians}

A common flow layer splits the variable into two blocks,
\[
z=(z_a,z_b),
\]
and transforms only one block at a time using the other as context.
An affine coupling layer has the form
\[
x_a = z_a,
\qquad
x_b = z_b \odot \exp\bigl(s(z_a)\bigr) + t(z_a),
\]
where $s(\cdot)$ and $t(\cdot)$ are arbitrary functions and $\odot$ denotes elementwise multiplication.

\begin{proposition}[Why coupling layers are tractable]
The Jacobian matrix of an affine coupling layer is block triangular. In particular,
\[
J =
\begin{pmatrix}
I & 0\\
\ast & \operatorname{diag}(\exp(s(z_a)))
\end{pmatrix},
\]
so that
\[
\det J = \exp\!\left(\sum_i s_i(z_a)\right).
\]
Moreover, the inverse map is explicit:
\[
z_a=x_a,
\qquad
z_b=\bigl(x_b-t(x_a)\bigr)\odot \exp\bigl(-s(x_a)\bigr).
\]
\end{proposition}

\begin{proof}
Because $x_a=z_a$, the derivative of $x_a$ with respect to $z_b$ is zero.
Thus the Jacobian is block lower triangular.
The diagonal block corresponding to $z_b\mapsto x_b$ is simply
\[
\operatorname{diag}(\exp(s(z_a))),
\]
whose determinant is the product of its diagonal entries.
The inverse formula follows by solving the affine transformation for $z_b$.
\end{proof}

\paragraph{Interpretation.}
Coupling layers are powerful because the functions $s$ and $t$ may be expressive neural networks, yet the determinant remains easy to compute.
This is the basic architectural trick that makes exact density modeling feasible in high dimensions.

\subsection*{A worked two-dimensional example}

Consider a two-dimensional standard Gaussian base variable
\[
Z=(Z_1,Z_2) \sim \mathcal{N}(0,I_2),
\]
and define the coupling transformation
\[
X_1 = Z_1,
\qquad
X_2 = e^{a Z_1} Z_2 + b Z_1,
\]
where $a,b\in\mathbb{R}$ are constants.
This is the special case
\[
s(z_1)=az_1, \qquad t(z_1)=bz_1.
\]

The inverse map is
\[
Z_1=X_1,
\qquad
Z_2=e^{-aX_1}(X_2-bX_1).
\]
The Jacobian matrix of the forward map is
\[
J_f(z_1,z_2)=
\begin{pmatrix}
1 & 0\\
aze^{az_1} z_2 + b & e^{az_1}
\end{pmatrix},
\]
so
\[
\det J_f(z_1,z_2)=e^{az_1}.
\]
Hence the model density is
\[
p_X(x_1,x_2)
=
\frac{1}{2\pi}
\exp\!\left(-\frac{1}{2}\Bigl[x_1^2+e^{-2ax_1}(x_2-bx_1)^2\Bigr]\right)
\cdot e^{-ax_1}.
\]

\paragraph{What the formula shows.}
Even though the base density is Gaussian, the transformed density need not be Gaussian.
The first coordinate controls both a translation and a scale change in the second coordinate.
Thus the map can bend, shear, and stretch the density while keeping exact likelihood available.

For example, if $a=0.5$ and $b=1$, then at the point $(x_1,x_2)=(1,2)$ we have
\[
z_1=1,
\qquad
z_2=e^{-0.5}(2-1)=e^{-0.5},
\]
and therefore
\[
\log p_X(1,2)
=
-\log(2\pi)
-\frac{1}{2}\left(1+e^{-1}\right)
-0.5.
\]
This gives a completely explicit likelihood calculation for a nontrivial nonlinear transport.

\subsection*{A geometric view: density warping}

Flows are often easiest to understand geometrically.
One begins with a simple cloud of probability mass in latent space and applies a sequence of invertible warpings.
Each layer preserves total probability but bends and stretches regions differently.
This can create complex curved level sets in data space while retaining exact density evaluation.

\begin{figure}[h]
\centering
\fbox{%
\begin{minipage}{0.88\textwidth}
\centering
\vspace{0.2em}
\textbf{Density warping under invertible transport}\par\medskip
\begin{tabular}{c c c}
Base space $z$ & $\xrightarrow{\quad f_1 \quad}$ & Intermediate space $h$ \\
concentric level sets &  & sheared / stretched level sets \\
\multicolumn{3}{c}{\vspace{0.4em}$\Downarrow$ composition of further invertible layers}\\[0.4em]
\multicolumn{3}{c}{Data space $x$: curved, anisotropic, non-Gaussian density contours}
\end{tabular}
\vspace{0.2em}
\end{minipage}}
\caption{A normalizing flow warps a simple base density into a complex target density through a sequence of invertible transformations. The important point is not only that samples are transported, but that local volume change remains exactly trackable.}
\end{figure}

\subsection*{What is elegant, and what is restrictive}

Flows appeal strongly to mathematically inclined readers because they provide an unusually clean synthesis of geometry and likelihood.
Still, this elegance comes with constraints.

\medskip

\noindent
\textbf{Strengths.}
\begin{itemize}
    \item exact likelihood evaluation,
    \item exact ancestral sampling,
    \item transparent relation between architecture and density geometry,
    \item direct maximum-likelihood training.
\end{itemize}

\noindent
\textbf{Limitations.}
\begin{itemize}
    \item invertibility constrains layer design,
    \item input and latent dimensions must match unless one augments the architecture,
    \item tractable Jacobians may restrict expressivity or efficiency,
    \item excellent likelihood does not always coincide with best perceptual sample quality.
\end{itemize}

Thus flows occupy a distinctive position in generative modeling: they are among the most exact and mathematically transparent models, but exactness imposes architectural discipline.

\subsection*{What to remember}

\begin{itemize}
    \item A normalizing flow models data by transporting a simple base density through an invertible map.
    \item The change-of-variables theorem gives exact density evaluation.
    \item Compositions turn global complexity into a sum of layerwise log-determinants.
    \item Coupling layers work because their Jacobians are triangular.
\end{itemize}

\subsection*{Common confusion}

\begin{itemize}
    \item Invertibility does not mean linearity; highly nonlinear flows are possible.
    \item Exact likelihood does not automatically imply the best visual samples.
    \item The determinant term is not an arbitrary correction; it is the precise account of local volume change.
\end{itemize}

\subsection*{Exercises}

\begin{enumerate}
    \item Verify directly that the inverse formula for the affine coupling layer is correct.
    \item Let $x=Az+b$ for an invertible matrix $A$. Derive the exact density of $x$ when $z\sim\mathcal{N}(0,I)$.
    \item For the two-dimensional example in this section, compute $p_X(x_1,x_2)$ when $a=0$ and explain why the model reduces to a simple shear transformation.
    \item Show that if each layer in a flow is volume preserving, then the log-determinant term vanishes and likelihood depends only on the base density evaluated at the inverse image.
\end{enumerate}

\subsection*{Notes and references}

The flow viewpoint goes back to the idea of exact density modeling through invertible transformations.
Important modern developments include NICE, RealNVP, Glow, and related architectures based on coupling layers, invertible $1\times 1$ convolutions, and other tractable Jacobian constructions.
Conceptually, flows also connect to transport maps and, at a broader level, to geometric perspectives on probability distributions.
